\documentclass[10pt,a4paper]{amsart}
\usepackage{amsmath,amssymb,amsthm,graphicx}
% Retained minimal preamble for the Stacks Project chapter sites-cohomology.tex (Cohomology on Sites).
% Verbatim excerpt of the chapter preamble: Xy-pic setup (native preamble lines 14-18) and the native
% theorem declarations with the mathematical macros (native preamble lines 154-217). Copied mechanically;
% nothing added, removed or redefined. The wrapper class and the amsmath/amssymb/amsthm/graphicx base
% packages are supplied by the assembler, so the class/externaldocument part of the native preamble is omitted.
% For commutative diagrams we use Xy-pic
\usepackage[all]{xy}

% We use 2cell for 2-commutative diagrams.
\xyoption{2cell}
% Theorem environments.
%
\theoremstyle{plain}
\newtheorem{theorem}[subsection]{Theorem}
\newtheorem{proposition}[subsection]{Proposition}
\newtheorem{lemma}[subsection]{Lemma}

\theoremstyle{definition}
\newtheorem{definition}[subsection]{Definition}
\newtheorem{example}[subsection]{Example}
\newtheorem{exercise}[subsection]{Exercise}
\newtheorem{situation}[subsection]{Situation}

\theoremstyle{remark}
\newtheorem{remark}[subsection]{Remark}
\newtheorem{remarks}[subsection]{Remarks}

\numberwithin{equation}{subsection}

% Macros
%
\def\lim{\mathop{\mathrm{lim}}\nolimits}
\def\colim{\mathop{\mathrm{colim}}\nolimits}
\def\Spec{\mathop{\mathrm{Spec}}}
\def\Hom{\mathop{\mathrm{Hom}}\nolimits}
\def\Ext{\mathop{\mathrm{Ext}}\nolimits}
\def\SheafHom{\mathop{\mathcal{H}\!\mathit{om}}\nolimits}
\def\SheafExt{\mathop{\mathcal{E}\!\mathit{xt}}\nolimits}
\def\Sch{\mathit{Sch}}
\def\Mor{\mathop{\mathrm{Mor}}\nolimits}
\def\Ob{\mathop{\mathrm{Ob}}\nolimits}
\def\Sh{\mathop{\mathit{Sh}}\nolimits}
\def\NL{\mathop{N\!L}\nolimits}
\def\CH{\mathop{\mathrm{CH}}\nolimits}
\def\proetale{{pro\text{-}\acute{e}tale}}
\def\etale{{\acute{e}tale}}
\def\QCoh{\mathit{QCoh}}
\def\Ker{\mathop{\mathrm{Ker}}}
\def\Im{\mathop{\mathrm{Im}}}
\def\Coker{\mathop{\mathrm{Coker}}}
\def\Coim{\mathop{\mathrm{Coim}}}

% Boxtimes
%
\DeclareMathSymbol{\boxtimes}{\mathbin}{AMSa}{"02}

%
% Macros for moduli stacks/spaces
%
\def\QCohstack{\mathcal{QC}\!\mathit{oh}}
\def\Cohstack{\mathcal{C}\!\mathit{oh}}
\def\Spacesstack{\mathcal{S}\!\mathit{paces}}
\def\Quotfunctor{\mathrm{Quot}}
\def\Hilbfunctor{\mathrm{Hilb}}
\def\Curvesstack{\mathcal{C}\!\mathit{urves}}
\def\Polarizedstack{\mathcal{P}\!\mathit{olarized}}
\def\Complexesstack{\mathcal{C}\!\mathit{omplexes}}
% \Pic is the operator that assigns to X its picard group, usage \Pic(X)
% \Picardstack_{X/B} denotes the Picard stack of X over B
% \Picardfunctor_{X/B} denotes the Picard functor of X over B
\def\Pic{\mathop{\mathrm{Pic}}\nolimits}
\def\Picardstack{\mathcal{P}\!\mathit{ic}}
\def\Picardfunctor{\mathrm{Pic}}
\def\Deformationcategory{\mathcal{D}\!\mathit{ef}}

% --- appended 2026-10-12 (Chief): the project own definitions of the reference/slogan/history
% display environments, copied mechanically from the shared Stacks preamble lines 8-12
% (\usepackage{verbatim} plus the three \newenvironment lines, which the chapter preamble keeps
% verbatim from the project preamble). Nothing added, removed or redefined: these environments
% carry no mathematical content and the project itself typesets them as comments. Without these
% three lines a counted statement whose own text uses them fails with Environment slogan undefined.
\usepackage{verbatim}
\newenvironment{reference}{\comment}{\endcomment}
\newenvironment{slogan}{\comment}{\endcomment}
\newenvironment{history}{\comment}{\endcomment}
\usepackage{cleveref}
\crefname{originallemma}{original lemma}{original lemmas}
\Crefname{originallemma}{Original lemma}{Original lemmas}
\makeatletter
\newlabel{categories-lemma-adjoint-fully-faithful}{{07RB}{source}{}{sourcepointer-1}{}}
\newlabel{categories-lemma-adjoint-fully-faithful@cref}{{[originallemma][0][]07RB}{[1][0][]source}{}{}{}}
\makeatother
\crefname{originalreference}{original reference}{original references}
\Crefname{originalreference}{Original reference}{Original references}
\makeatletter
\newlabel{example-category-to-point}{{08PF}{source}{}{sourcepointer-2}{}}
\newlabel{example-category-to-point@cref}{{[originalreference][0][]08PF}{[1][0][]source}{}{}{}}
\makeatother
\title{The lower shriek along the projection to the point is a left inverse of the pullback when the site has an initial or a final object}
\author{The Stacks Project Authors}
\date{stacks-project:sites-cohomology.tex:local-snapshot-20261011; GNU FDL 1.2 or later}
\begin{document}
\maketitle
\noindent\textit{Editorial note.} All delivered mathematical text is the author's own, copied from the pinned Stacks Project chapter file sites-cohomology.tex; no statement or proof step was rewritten, reordered or removed.\par\medskip
\noindent\textit{Editorial note.} The counted proof is the author's complete argument: the evaluation of the lower shriek on an initial object, the evaluation of the pushforward on a final object giving Rpi\_* pi\textasciicircum{}\{-1\} = id, and the formal conclusion through the criterion for full faithfulness of an adjoint.\par\medskip
\noindent\textit{Editorial note.} The label lemma-initial-final names the counted statement, and the proof environment immediately following it in the source is delivered verbatim as its proof.\par\medskip
\noindent\textit{Editorial note.} The counted statement is carried by the author's own lemma environment in the chapter source.\par\medskip
\noindent\textit{Editorial note.} Pre-open gate: python operations/agent-authoring/tools/precard\_gate.py --source sites-cohomology.tex --label lemma-initial-final --start 10188 --end 10206 --local-id 3570 returned RESULT=OK (L3570 free, occupancy FREE, public interval clear). Candidate taken from operations/pool-cards/stacks-sites-cohomology-halfB.json.\par\medskip
\noindent\textit{Editorial note.} The statement and proof are copied from the pinned chapter file; the cited inputs are kept as external pointers with their canonical Stacks tags.\par\medskip
\noindent\textit{Allowed standard prerequisites (editorial).}\begin{itemize}
\item Computation of the lower shriek and of the pushforward on a site with an initial or final object. Same chapter, the example computing the functors pi\_! and Rpi\_* for the projection to the point when the site has an initial or a final object. Supplies the evaluation of pi\_! on an initial object and of Rpi\_* on a final object used in the two cases.
\item Full faithfulness of an adjoint whose unit or counit is an isomorphism. Categories, the lemma that a functor with a left adjoint whose counit is an isomorphism is fully faithful and the subsequent statement identifying the adjunction when both maps are isomorphisms. Supplies the conclusion that pi\textasciicircum{}\{-1\} is fully faithful when Rpi\_* pi\textasciicircum{}\{-1\} = id, and then that the counit pi\_! pi\textasciicircum{}\{-1\} -> id is an isomorphism.
\end{itemize}\medskip
\ifdefined\hypertarget\hypertarget{sourcepointer-1}{}\fi
\noindent\textit{Source reference.} Original lemma 07RB: sources/stacks-chow/tags.txt, source lines 9861--9861. Prerequisite pointer invoked at proof line(s) 10203; it names a standard fact of a sibling chapter that the counted argument uses, and every step of the counted proof built on it is retained verbatim in the delivered body.\par\medskip
\ifdefined\hypertarget\hypertarget{sourcepointer-2}{}\fi
\noindent\textit{Source reference.} Original reference 08PF: sources/stacks-sites-cohomology/pinned-original-source.tex, source lines 10016--10016. Prerequisite pointer invoked at proof line(s) 10190; it names a standard fact of the same chapter that the counted argument uses, and every step of the counted proof built on it is retained verbatim in the delivered body.\par\medskip
\begin{lemma}
\label{lemma-initial-final}
Notation and assumptions as in Example \ref{example-category-to-point}.
If $\mathcal{C}$ has either an initial or a final object, then
$L\pi_! \circ \pi^{-1} = \text{id}$ on $D(\textit{Ab})$, resp.\ $D(B)$.
\end{lemma}

\begin{proof}
If $\mathcal{C}$ has an initial object, then $\pi_!$ is computed by
evaluating on this object and the statement is clear. If $\mathcal{C}$
has a final object, then $R\pi_*$ is computed by evaluating on this
object, hence $R\pi_* \circ \pi^{-1} \cong \text{id}$ on
$D(\textit{Ab})$, resp.\ $D(B)$. This implies that
$\pi^{-1} : D(\textit{Ab}) \to D(\mathcal{C})$,
resp.\ $\pi^{-1} : D(B) \to D(\underline{B})$ is fully faithful, see
Categories, Lemma \ref{categories-lemma-adjoint-fully-faithful}.
Then the same lemma implies that $L\pi_! \circ \pi^{-1} = \text{id}$
as desired.
\end{proof}

\end{document}
